% This latex document tries to copy the microsoft word template of XXXIV CILAMCE

\documentclass[a4paper,12pt,fleqn]{article}

\usepackage{graphicx}
\usepackage{amsmath,amsfonts, amsthm}
\usepackage{enumerate}
\usepackage[shortlabels]{enumitem} 
\usepackage{tabularx}
\usepackage{comment}
\usepackage{url}
\usepackage{color}
\usepackage{cancel}
\usepackage{subfigure}
\usepackage{xcolor}



% PACKAGES USED - packages that need to be previously installed on your computer
\usepackage[lmargin=2.5cm, rmargin=2.5cm, tmargin=2.5cm, bmargin=3.5cm ]{geometry}
\usepackage{graphicx}
%\usepackage{times}
\usepackage{indentfirst}
\usepackage{fancyhdr}
%\usepackage{titlesec}
%\usepackage[english]{babel}      % nomes em português
%\selectlanguage{portuges}
\usepackage[english,brazil]{babel} %Opções de escrita em inglês, português para Figura, Tabela.
\usepackage{parskip}
\usepackage{setspace}

%%% Overleaf:
\usepackage[utf8]{inputenc}
%%% Texmaker:
% \usepackage[latin1]{inputenc}
\usepackage[brazil]{babel}


\usepackage{amssymb} % Usado para desigualdades nas equações
\usepackage{dsfont} %Usado para conjuntos N, Z, Q, R, C


\usepackage{cite}


\usepackage[pdftex]{hyperref}    % "links" dentro do pdf
\usepackage[T1]{fontenc}          % permitir \hyphenation

\usepackage{pifont}				% caracteres decorativos
%\usepackage[utf8]{inputenc} % Usado para acentuações



% CONFIGURATION
\renewcommand*\arraystretch{1.5} % height of the table's row
\usepackage[font=bf, size=footnotesize]{caption} % bold face font for labels of figures and tables
\renewcommand*\thesection{\arabic{section}} % for use book class instead article class
%\hyphenpenalty=10000 % uncomment for avoir hypenization
%\titleformat*{\section}{\large\bfseries}
%\titleformat*{\subsection}{\large\bfseries}
%\titlespacing\section{0pt}{12pt plus 12pt minus 6pt}{6pt plus 6pt minus 6pt}
%\titlespacing\subsection{0pt}{6pt plus 6pt minus 6pt}{6pt plus 6pt minus 6pt}
%\titlespacing\subsubsection{0pt}{6pt plus 6pt minus 6pt}{6pt plus 6pt minus 6pt}
\setlength{\parskip}{6pt}
\setlength{\parindent}{0.75cm}
\let\ds=\displaystyle
% PDF zoom {the last parameter is the zoom: 1.00 (100%), 0.5 (50%), ...}
\usepackage{hyperref}
\hypersetup{pdfstartview={XYZ null null 1.00}}

\usepackage{lipsum}

% Para sistemas lineares:
\usepackage[overload]{empheq}

% Linhas tracejadas 
\usepackage{arydshln}
\usepackage{mathtools}


\begin{document}

\begin{center}

{\color{blue}

\large
Cálculo Numérico - IME/UERJ 

Gabarito - Lista de Exercícios 3 - Sistemas Lineares - Métodos diretos

}
\end{center}

\vspace{0.8cm}

  
  \begin{enumerate}
  
  \item 
  
      \begin{enumerate}

         {\color{blue} 
         
         \item Soluções: $ \begin{bmatrix}  -0,0000990 & 0,0000098 \end{bmatrix}^T $; $ \begin{bmatrix}  0,0098029 & 0,0990294 \end{bmatrix}^T $
         
         
         \item Mau condicionamento da matriz (Por quê?)
         
         }
         
       \end{enumerate}
       
  
  \item
  
  {\color{blue}
  
  Conjunto de soluções do sistema linear é infinito e está definido por:
     
     $S = \Big\{ (x,y,z) \in \mathds{R}^3 ~\Big \vert ~ x = \dfrac{1}{16}(z+13) ; y = \dfrac{1}{8} (11z - 17) \Big\} $
     
  }
  
  \item 
     
  {\color{blue} Há necessidade de troca da primeira e terceira linhas. 
     
     $(x, y, z) = (10/7, -5/3, 9/5) \approx (1,4286; -1,6667; 1,8000)$, usando 4 dígitos e arredondamento.
     
  }
     
     
%  \item
%  
%  \begin{enumerate}[(a)]
% \item 
%
%\begin{enumerate}[(i)]
%\item  {\color{blue} $X = (6/5, -1/5)^t$;  }
%\item {\color{blue} $X = (-3/5, -2/5)^t$;  }
%\end{enumerate}
%
% \item 
% {\color{blue}
% 
% \[
%  A^{-1} =
%  \left[ {\begin{array}{rr}
%    1/5 & 4/5 \\
%    -1/5 &  1/5 \\
%  \end{array} } \right].
%\]
% 
%}
% 
%\end{enumerate}
  
\vspace{0.3cm} 
    
 \item 
 

\begin{enumerate}

 {\color{blue} 


\item  

\[ U =
    \left[ {\begin{array}{rrr}
		2 & -2 & 3 \\
		{\color{red} 0 } & 6 & -5 \\
		{\color{red} 0 } & {\color{red} 0} & 13/3 
  \end{array} } \right];
\]


\[ L =
    \left[ {\begin{array}{rrr}
		1 & 0 & 0 \\
		{\color{red} m_{21} } & 1 & 0 \\
		{\color{red} m_{31} } & {\color{red} m_{32}} & 1 
  \end{array} } \right]
  =
      \left[ {\begin{array}{rrr}
		1 & 0 & 0 \\
		{\color{red} 2 } & 1 & 0 \\
		{\color{red} -3 } & {\color{red} -4/3} & 1 
  \end{array} } \right]
\]

\item

Devemos resolver o sistema linear $Ax = b$ com as matrizes $L$ e $U$ encontradas.

No caso sem pivoteamento, temos $A = LU$. Assim,
     
     $Ax = b \Rightarrow LUx = b \Rightarrow \left\{\begin{array}{rcl} Ly &=& b \quad (1) \Rightarrow \text{Acho } y \\ Ux &=& y \quad (2) ~\Rightarrow \text{Acho } x  \end{array}\right. $

\[ x =
  \left[ \begin{array}{r}
		x_1\\
		x_2\\
		x_3
  \end{array} \right]
  =
   \left[ \begin{array}{r}
		0 \\
		0 \\
		1
  \end{array} \right].
\]


 }    
 
\end{enumerate}
  

\item 


\begin{enumerate}


 {\color{blue} 
 
\item
 
Devemos achar as matrizes $L$, $U$ e $P$ usando eliminação de Gauss \textbf{com pivoteamento parcial}. 

\[ U =
    \left[ {\begin{array}{rrr}
		-4 & -1 & 2 \\
		{\color{red} 0 } & 1 & 1 \\
		{\color{red} 0 } & {\color{red} 0} & 15/4 
  \end{array} } \right];
\]

\[ L =
    \left[ {\begin{array}{rrr}
		1 & 0 & 0 \\
		{\color{red} 0 } & 1 & 0 \\
		{\color{red} -1/4 } & {\color{red} 3/4} & 1 
  \end{array} } \right]; 
\]

\[ P = P^{(1)} P^{(0)} = 
    \left[ {\begin{array}{rrr}
		1 & 0 & 0 \\
		0 & 0 & 1 \\
		0 & 1 & 0 
  \end{array} } \right]
    \left[ {\begin{array}{rrr}
		0 & 0 & 1 \\
		0 & 1 & 0 \\
		1 & 0 & 0 
  \end{array} } \right]
  =
    \left[ {\begin{array}{rrr}
		0 & 0 & 1 \\
		1 & 0 & 0 \\
		0 & 1 & 0 
  \end{array} } \right].
\]


\item 

PARTE 2: Devemos resolver o sistema linear $Ax = b$ com as matrizes $L$, $U$ e $P$ encontradas.

No caso com pivoteamento, temos $PA = LU$. Assim,
     
     $Ax = b \Rightarrow PAx = Pb \Rightarrow LUx = Pb \Rightarrow \left\{\begin{array}{rcl} Ly &=& Pb \quad (1) \Rightarrow \text{Acho } y \\ Ux &=& y \quad (2) ~\Rightarrow \text{Acho } x  \end{array}\right. $
     


\[ x =
  \left[ \begin{array}{r}
		x_1\\
		x_2\\
		x_3
  \end{array} \right]
  =
   \left[ \begin{array}{r}
		6/5 \\
		-8/5 \\
		3/5
  \end{array} \right].
\]


 }    
 
  
\end{enumerate}


\item 

{\color{blue}

    \begin{enumerate}
    
     \item $A = LU \Rightarrow A^{-1} = U^{-1} L^{-1}$.
     
     Primeiro, achamos $L^{-1}$ por eliminação de Gauss na transformação
     \begin{equation*}
   		  \setlength{\dashlinegap}{2pt}
  		  \left[
  		  		\begin{array}{c:c}
				   A & I 
		        \end{array}
		  \right]
		  \Rightarrow
   		  \setlength{\dashlinegap}{2pt}
  		  \left[
  		  		\begin{array}{c:c}
				   U & L ^{-1} 
		        \end{array}
		  \right]
	\end{equation*}
	
	de onde tiramos 
	\[ L^{-1} =
    \left[ {\begin{array}{rrr}
		1 & 0 & 0 \\
		{\color{red} -2 } & 1 & 0 \\
		{\color{red} 1/2 } & {\color{red} 3/2} & 1 
  \end{array} } \right]; 
   \]
	
	Depois, achamos achamos $U^{-1}$ por Gauss-Jordan na transformação
     \begin{equation*}
   		  \setlength{\dashlinegap}{2pt}
  		  \left[
  		  		\begin{array}{c:c}
				   U & I 
		        \end{array}
		  \right]
		  \Rightarrow
   		  \setlength{\dashlinegap}{2pt}
  		  \left[
  		  		\begin{array}{c:c}
				   I & U ^{-1} 
		        \end{array}
		  \right]
	\end{equation*}

    de onde tiramos 
	\[ U^{-1} =
    \left[ {\begin{array}{rrr}
		1/2 & 1/6 & -2/15 \\
		{\color{red} 0 } & 1/6 & 1/6 \\
		{\color{red} 0 } & {\color{red} 0} & 1/5 
  \end{array} } \right]; 
   \]
   
   Assim,
   \[ A^{-1} =
    \left[ {\begin{array}{rrr}
		1/2 & 1/6 & -2/15 \\
		{\color{red} 0 } & 1/6 & 1/6 \\
		{\color{red} 0 } & {\color{red} 0} & 1/5 
  \end{array} } \right]
    \left[ {\begin{array}{rrr}
		1 & 0 & 0 \\
		{\color{red} -2 } & 1 & 0 \\
		{\color{red} 1/2 } & {\color{red} 3/2} & 1 
  \end{array} } \right].
   \]
   
   
   
   \item Com pivoteamento parcial: $PA = LU \Rightarrow A^{-1} = U^{-1} L^{-1} P$, onde $P$ é a matriz das permutações de linhas na matriz identidade.
   
   Primeiro, achamos $L^{-1}P$ por eliminação de Gauss com pivoteamento parcial na transformação
    \begin{equation*}
   		  \setlength{\dashlinegap}{2pt}
  		  \left[
  		  		\begin{array}{c:c}
				   A & I 
		        \end{array}
		  \right]
		  \Rightarrow
   		  \setlength{\dashlinegap}{2pt}
  		  \left[
  		  		\begin{array}{c:c}
				   U & L ^{-1} P 
		        \end{array}
		  \right]
	\end{equation*}  
	
	Depois, achamos achamos $U^{-1}$ por Gauss-Jordan na transformação
	\begin{equation*}
   		  \setlength{\dashlinegap}{2pt}
  		  \left[
  		  		\begin{array}{c:c}
				   U & I 
		        \end{array}
		  \right]
		  \Rightarrow
   		  \setlength{\dashlinegap}{2pt}
  		  \left[
  		  		\begin{array}{c:c}
				   I & U ^{-1} 
		        \end{array}
		  \right]
	\end{equation*}   
  
    (Faça as contas!)
  
%\begin{equation*}
%\setlength{\dashlinegap}{2pt}
%\left[\begin{array}{cccc:c}
%a_{11} & a_{12} & \cdots & a_{1n} & b_1 \\
%a_{21} & a_{22} & \cdots & a_{2n} & b_2 \\
%\vdots & \vdots & \ddots & \vdots & \vdots \\
%a_{n1} & a_{n2} & \cdots & a_{nn} & b_n
%\end{array}
%\right]
%\end{equation*}
    
    
    \end{enumerate}
    
}
 

%\item
%
%\begin{enumerate}[(a)]
%     
%\item  {\color{blue} $X = (0.3906,1.2031);$ }
%\item  {\color{blue} $X = (0.4023,1.2012);$ }
%
%\end{enumerate}
%
%\item           
%
%      \begin{enumerate}[(a)]
%        \item 
%        
%        {\color{blue}
%			
%			\textbf{Resposta:} Devemos trocar a primeira e a terceira equações (ou a primeira e a terceira linhas da matriz $A$ e do vetor $b$). Assim, obtemos:
%			
%\[
%  A' =
%  \left[ {\begin{array}{rrr}
%		4 & ~~~1 & ~~~2 \\
%		3 & ~-5 & ~~~0 \\
%		2 & ~-4 & ~-7
%  \end{array} } \right];
%  \qquad \qquad b' =   \left[ {\begin{array}{rrrr}
%    1 \\
%    -5 \\
%    0
%  \end{array} } \right].
%\]			
%			Usando o Critério de Sassenfeld, obtemos $\beta = \max \left \{ \ds \frac{3}{4}, \ds \frac{9}{20}, \ds \frac{33}{70} \right \} = \ds \frac{3}{4} < 1 $.
%			
%			Logo, com essa reordenação de equações, há garantia de convergência usando o método iterativo de Gauss-Seidel.
%
%        }
%        
%
%       
%        \item 
%        
%              {\color{blue}
%			
%			     \textbf{Resposta:} $X^{(1)} = (0,2575;1,1545;-0,5861)^T$.
%			
%			  }
%        
%        \item
%        
%              {\color{blue}
%			
%			     \textbf{Resposta:} $4,5 \times 10^{-3}$.
%			
%			  }
%        
%      \end{enumerate}     
%      
%\vspace{0.3cm} 
%
%\item    
%
%{\color{blue}
%
%\textbf{Resposta:}
%
%Primeiro, verifique que o Critério de Sassenfeld é satisfeito.
%
%Resolvendo o sistema pelo método de Gauss-Seidel, uma primeira aproximação $X^{(3)}$, partindo de $X^{(0)} = (0,0,0)^t$, é dada por:
%
%$X^{(3)} = (0.7440,-0.7381,2.0952)^t$.
%
%}
%
%\vspace{0.3cm} 
%
%      
%\item           
%      
%      \begin{enumerate}[(a)]
%        \item 
%        {\color{blue} 
%        Sim. 
%        }
%        
%        \item 
%        {\color{blue} 
%        $X = (0.3636,0.4545,0.4545,0.3636)^t$ 
%        }      
%      \end{enumerate} \vspace{0.3cm}          
      
   


%\vspace{0.3cm}
%
%\item           
%      
%      \begin{enumerate}[(a)]
%        \item 
%        {\color{blue} 
%        $\alpha_1 = \dfrac{ \vert 2 \vert + \vert -1 \vert + \vert 0 \vert}{\vert 1 \vert} = \dfrac{2}{2.1} = {\color{red} 3 > 1  \Rightarrow \text{Não satisfaz o Critério das Linhas.} }$ 
%        }
%        
%        \item 
%        {\color{blue} 
%        $\beta_1 = \alpha_1 = {\color{red} 3 > 1  \Rightarrow \text{Não satisfaz o Critério de Sassenfeld.} }$ 
%        }   
%        
%        \item 
%        {\color{blue} 
%        Resolvendo o sistema a partir de  $X^{(0)} = (0,0,0,0)^t$, não há convergência em nenhum dos métodos.
%        }           
%        
%        \item 
%        {\color{blue} 
%
%        Permutando-se as duas primeiras equações para o sistema, o Critério de Sassenfeld é satisfeito (verifique).
%        
%%Permutando-se as duas primeiras equações, obtemos o sistema:
%%
%%\begin{alignat*}{6}[left = \empheqlbrace]
%%    \quad 2x_1 & \quad - & \quad x_2 & {}{} &  &  & & & \quad = &  \quad 1\phantom{0} \\
%%   \quad x_1 & \quad + & \quad 2x_2 & \quad - & \quad x_3 &  & & &\quad = & \quad 1 \\
%%     & \quad - & \quad x_2 & \quad + & \quad 2x_3 & \quad - & \quad x_4 & & \quad = & \quad 1 \\
%%     &  &  & \quad - & \quad x_3 & \quad + & \quad 2x_4 & &\quad = & \quad 1 
%%\end{alignat*}  
%%
%%onde:
%%
%%$\beta_1 = \dfrac{ \vert - 1 \vert + \vert 0 \vert + \vert 0 \vert}{\vert 2 \vert} = 0.5 < 1$
%%
%%$\beta_2 = \dfrac{ 0.5\vert 1 \vert + \vert -1 \vert + \vert 0 \vert}{\vert 2 \vert} = 0.75 < 1$
%%
%%$\beta_3 = \dfrac{ 0.5\vert 0 \vert + 0.75 \vert -1 \vert + \vert -1 \vert}{\vert 2 \vert} = 0.875 < 1$
%%
%%$\beta_4 = \dfrac{ 0.5\vert 0 \vert + 0.75 \vert 0 \vert + 0.875 \vert -1 \vert}{\vert 2 \vert} = 0.4375 < 1$
%%
%%$\max\{\beta_1, \beta_2, \beta_3, \beta_4\} = 0.875 <  1$
%%
%%Logo, o Critério de Sassenfeld é satisfeito.
%
%}
%
%        \item 
%        {\color{blue} 
%        Resolvendo o sistema do item (d) a partir de  $X^{(0)} = (0,0,0,0)^t$, obtemos a solução:
%        
%        $X^{*} = (0.9052, ~0.8150, ~1.5413, ~1.2706)^t$.
%		        
%        
%        }           
%        
%      \end{enumerate} 

      
\end{enumerate}


\end{document}
